\chapter{Improper Integrals}

An improper integral is one where either
the interval of integration is unbounded, or the integrand has a
vertical asymptote inside the interval. Both are defined as limits of
proper integrals and may or may not converge. The economist's archetype
is $\int_{0}^{\infty} \mathrm{e}^{-rt}\,\mathrm{d}t$, the present value
of a unit perpetual flow.

\section{Core Rules}

\begin{rulebox}{Type I: unbounded interval}
For continuous $f$,
\[
  \int_{a}^{\infty} f(x)\,\mathrm{d}x \;:=\; \lim_{b\to\infty}\int_{a}^{b} f(x)\,\mathrm{d}x,
\]
\[
  \int_{-\infty}^{b} f(x)\,\mathrm{d}x \;:=\; \lim_{a\to-\infty}\int_{a}^{b} f(x)\,\mathrm{d}x,
\]
\[
  \int_{-\infty}^{\infty} f(x)\,\mathrm{d}x \;:=\; \int_{-\infty}^{c} f(x)\,\mathrm{d}x+\int_{c}^{\infty} f(x)\,\mathrm{d}x,
\]
the last requiring \emph{both} pieces to converge separately
(independent of the splitting point $c$).
\end{rulebox}

\begin{rulebox}{Type II: vertical asymptote}
If $f$ is continuous on $[a,b)$ and blows up at $b$,
\[
  \int_{a}^{b} f(x)\,\mathrm{d}x \;:=\; \lim_{t\to b^{-}}\int_{a}^{t} f(x)\,\mathrm{d}x .
\]
Analogously for an asymptote at $a$ (limit from $a^{+}$), or at an
interior point $c$ (split at $c$ and treat each piece).
\end{rulebox}

\begin{rulebox}{Convergence vs.\ divergence}
If the limit exists and is finite, the improper integral
\emph{converges} to that value; otherwise it \emph{diverges} (either to
$\pm\infty$, or by oscillation).
\end{rulebox}

\begin{rulebox}{Two benchmark families}
\begin{itemize}[leftmargin=1.6em,itemsep=0.3ex]
  \item $\displaystyle\int_{1}^{\infty} \frac{1}{x^{p}}\,\mathrm{d}x$
        converges iff $p>1$, and equals $\dfrac{1}{p-1}$ in that case.
  \item $\displaystyle\int_{0}^{1} \frac{1}{x^{p}}\,\mathrm{d}x$
        converges iff $p<1$, and equals $\dfrac{1}{1-p}$ in that case.
\end{itemize}
The borderline case $p=1$ diverges in both.
\end{rulebox}

\begin{rulebox}{The economist's identity}
For any $r>0$,
\[
  \int_{0}^{\infty} \mathrm{e}^{-rt}\,\mathrm{d}t=\frac{1}{r},\qquad
  \int_{0}^{\infty} t\,\mathrm{e}^{-rt}\,\mathrm{d}t=\frac{1}{r^{2}},\qquad
  \int_{0}^{\infty} \mathrm{e}^{-rt}f(t)\,\mathrm{d}t=\text{PV of cash flow }f .
\]
\end{rulebox}

\section{Practice Problems}

Decide whether each integral converges; if it does, evaluate it.

\begin{problems}
  \item $\displaystyle\int_{1}^{\infty} \frac{1}{x^{2}}\,\mathrm{d}x$
  \item $\displaystyle\int_{1}^{\infty} \frac{1}{x}\,\mathrm{d}x$
  \item $\displaystyle\int_{1}^{\infty} \frac{1}{\sqrt{x}}\,\mathrm{d}x$
  \item $\displaystyle\int_{0}^{\infty} \mathrm{e}^{-x}\,\mathrm{d}x$
  \item $\displaystyle\int_{0}^{\infty} \mathrm{e}^{-rt}\,\mathrm{d}t$
        \quad ($r>0$)
  \item $\displaystyle\int_{0}^{\infty} t\,\mathrm{e}^{-t}\,\mathrm{d}t$
  \item $\displaystyle\int_{0}^{\infty} t^{2}\,\mathrm{e}^{-t}\,\mathrm{d}t$
  \item $\displaystyle\int_{1}^{\infty} \frac{1}{x^{3}}\,\mathrm{d}x$
  \item $\displaystyle\int_{-\infty}^{0} \mathrm{e}^{x}\,\mathrm{d}x$
  \item $\displaystyle\int_{-\infty}^{\infty} \frac{x}{1+x^{2}}\,\mathrm{d}x$
  \item $\displaystyle\int_{0}^{1} \frac{1}{\sqrt{x}}\,\mathrm{d}x$
  \item $\displaystyle\int_{0}^{1} \frac{1}{x}\,\mathrm{d}x$
  \item $\displaystyle\int_{0}^{1} \frac{1}{x^{1/3}}\,\mathrm{d}x$
  \item $\displaystyle\int_{0}^{1} \ln x\,\mathrm{d}x$
  \item $\displaystyle\int_{1}^{\infty} \frac{\ln x}{x^{2}}\,\mathrm{d}x$
  \item $\displaystyle\int_{2}^{\infty} \frac{1}{x(\ln x)^{2}}\,\mathrm{d}x$
  \item $\displaystyle\int_{2}^{\infty} \frac{1}{x\ln x}\,\mathrm{d}x$
  \item $\displaystyle\int_{0}^{\infty} \frac{1}{(1+x)^{2}}\,\mathrm{d}x$
  \item $\displaystyle\int_{-\infty}^{\infty} \frac{1}{1+x^{2}}\,\mathrm{d}x$
        \quad (antideriv.\ is $\arctan x$)
  \item Present value of a perpetuity that pays continuously at rate
        $D=1000$ per year, discounted at $r=0.05$:
        $\displaystyle\int_{0}^{\infty} 1000\,\mathrm{e}^{-0.05t}\,\mathrm{d}t$.
  \item Present value of a finite annuity:
        $\displaystyle\int_{0}^{T} D\,\mathrm{e}^{-rt}\,\mathrm{d}t$.
  \item $\displaystyle\int_{0}^{2} \frac{1}{(x-2)^{2}}\,\mathrm{d}x$
        \quad (asymptote at $x=2$)
  \item $\displaystyle\int_{0}^{1} \frac{1}{(1-x)^{1/2}}\,\mathrm{d}x$
  \item $\displaystyle\int_{0}^{\infty} x\,\mathrm{e}^{-x^{2}}\,\mathrm{d}x$
  \item $\displaystyle\int_{1}^{\infty} \frac{1}{x^{2}-x}\,\mathrm{d}x
        = \int_{1}^{\infty} \frac{1}{x(x-1)}\,\mathrm{d}x$
        \quad (careful! asymptote at $x=1$)
\end{problems}

\section{Solutions}

\begin{solutions}
  \item $\displaystyle\Bigl[-\frac{1}{x}\Bigr]_{1}^{b}=-\frac{1}{b}+1\to 1$
        as $b\to\infty$. Converges to $1$.

  \item $\displaystyle\bigl[\ln x\bigr]_{1}^{b}=\ln b\to\infty$. Diverges.

  \item $\displaystyle\bigl[2\sqrt{x}\bigr]_{1}^{b}=2\sqrt{b}-2\to\infty$.
        Diverges.

  \item $\displaystyle\bigl[-\mathrm{e}^{-x}\bigr]_{0}^{b}=1-\mathrm{e}^{-b}\to 1$.
        Converges to $1$.

  \item $\displaystyle\Bigl[-\frac{1}{r}\mathrm{e}^{-rt}\Bigr]_{0}^{b}
        =\frac{1}{r}(1-\mathrm{e}^{-rb})\to\frac{1}{r}$. Converges to
        $1/r$. This is the present value of a unit perpetual flow.

  \item By parts (Section 3 problem 11 with $r=1$): an antiderivative is
        $-(t+1)\mathrm{e}^{-t}$.
        \[
          \bigl[-(t+1)\mathrm{e}^{-t}\bigr]_{0}^{b}=-(b+1)\mathrm{e}^{-b}+1\to 1 .
        \]
        Converges to $1$. (Generalised:
        $\int_{0}^{\infty} t\mathrm{e}^{-rt}\,\mathrm{d}t=1/r^{2}$.)

  \item Tabular IBP gives antiderivative
        $-(t^{2}+2t+2)\mathrm{e}^{-t}$ (Section 3 problem 13).
        \[
          \bigl[-(t^{2}+2t+2)\mathrm{e}^{-t}\bigr]_{0}^{b}
            =-(b^{2}+2b+2)\mathrm{e}^{-b}+2\to 2 .
        \]
        Converges to $2$. (This is $\Gamma(3)=2!$.)

  \item $\displaystyle\Bigl[-\frac{1}{2x^{2}}\Bigr]_{1}^{b}
        =-\frac{1}{2b^{2}}+\frac{1}{2}\to\frac{1}{2}$. Converges to $1/2$.

  \item $\displaystyle\bigl[\mathrm{e}^{x}\bigr]_{a}^{0}=1-\mathrm{e}^{a}\to 1$
        as $a\to-\infty$. Converges to $1$.

  \item Split at $0$.
        $\int_{0}^{\infty} \frac{x}{1+x^{2}}\,\mathrm{d}x
        =\lim_{b}\frac{1}{2}\ln(1+b^{2})=\infty$. The right piece
        already diverges, so the doubly improper integral
        \textbf{diverges} (even though the integrand is odd: the
        principal value is $0$ but the standard definition requires
        both halves to converge separately).

  \item $\displaystyle\bigl[2\sqrt{x}\bigr]_{t}^{1}=2-2\sqrt{t}\to 2$
        as $t\to 0^{+}$. Converges to $2$.

  \item $\displaystyle\bigl[\ln x\bigr]_{t}^{1}=-\ln t\to\infty$ as
        $t\to 0^{+}$. Diverges.

  \item $\displaystyle\int_{t}^{1} x^{-1/3}\,\mathrm{d}x
        =\Bigl[\tfrac{3}{2}x^{2/3}\Bigr]_{t}^{1}
        =\tfrac{3}{2}(1-t^{2/3})\to\tfrac{3}{2}$. Converges to $3/2$.

  \item Antiderivative $x\ln x-x$ (Section 3 problem 4).
        \[
          \bigl[x\ln x-x\bigr]_{t}^{1}=(0-1)-(t\ln t-t)=-1-t\ln t+t .
        \]
        As $t\to 0^{+}$, $t\ln t\to 0$, so the limit is $-1$. Converges
        to $-1$.

  \item By parts: $u=\ln x$, $v'=x^{-2}$, $v=-x^{-1}$.
        \[
          \int \frac{\ln x}{x^{2}}\,\mathrm{d}x
            =-\frac{\ln x}{x}+\int \frac{1}{x^{2}}\,\mathrm{d}x
            =-\frac{\ln x}{x}-\frac{1}{x}+C=-\frac{\ln x+1}{x}+C .
        \]
        Evaluate $\displaystyle\Bigl[-\frac{\ln x+1}{x}\Bigr]_{1}^{b}
        =-\frac{\ln b+1}{b}+1\to 1$ (numerator grows like $\ln b$,
        denominator like $b$). Converges to $1$.

  \item $u=\ln x$, $\mathrm{d}u=\mathrm{d}x/x$. Antiderivative
        $-\dfrac{1}{\ln x}$.
        \[
          \Bigl[-\frac{1}{\ln x}\Bigr]_{2}^{b}
            =-\frac{1}{\ln b}+\frac{1}{\ln 2}\to\frac{1}{\ln 2} .
        \]
        Converges to $1/\ln 2$.

  \item Same substitution; antiderivative $\ln|\ln x|$.
        \[
          \bigl[\ln\ln x\bigr]_{2}^{b}=\ln\ln b-\ln\ln 2\to\infty .
        \]
        Diverges (just barely).

  \item $u=1+x$. Antiderivative $-\dfrac{1}{1+x}$.
        \[
          \Bigl[-\frac{1}{1+x}\Bigr]_{0}^{b}
            =-\frac{1}{1+b}+1\to 1 .
        \]
        Converges to $1$.

  \item $\displaystyle\bigl[\arctan x\bigr]_{-a}^{b}
        \to\frac{\pi}{2}-\bigl(-\frac{\pi}{2}\bigr)=\pi$. Converges to
        $\pi$.

  \item By problem 5 with $r=0.05$:
        \[
          PV=\frac{1000}{0.05}=20{,}000 .
        \]

  \item $\displaystyle\int_{0}^{T} D\,\mathrm{e}^{-rt}\,\mathrm{d}t
        =D\bigl[-\tfrac{1}{r}\mathrm{e}^{-rt}\bigr]_{0}^{T}
        =\frac{D(1-\mathrm{e}^{-rT})}{r}$. As $T\to\infty$ this tends
        to $D/r$, recovering the perpetuity.

  \item Antiderivative $-\dfrac{1}{x-2}$. As $t\to 2^{-}$,
        \[
          \Bigl[-\frac{1}{x-2}\Bigr]_{0}^{t}
            =-\frac{1}{t-2}-\frac{1}{2}=\frac{1}{2-t}-\frac{1}{2}\to\infty .
        \]
        Diverges.

  \item $u=1-x$, $\mathrm{d}u=-\mathrm{d}x$. Antiderivative
        $-2\sqrt{1-x}$.
        \[
          \bigl[-2\sqrt{1-x}\bigr]_{0}^{t}=-2\sqrt{1-t}+2\to 2 .
        \]
        Converges to $2$.

  \item $u=x^{2}$, $\mathrm{d}u=2x\,\mathrm{d}x$. Antiderivative
        $-\tfrac12\mathrm{e}^{-x^{2}}$.
        \[
          \Bigl[-\tfrac12\mathrm{e}^{-x^{2}}\Bigr]_{0}^{b}
            =-\tfrac12\mathrm{e}^{-b^{2}}+\tfrac12\to\tfrac12 .
        \]
        Converges to $1/2$.

  \item The integrand also blows up at the left endpoint $x=1$, so the
        integral is improper at \emph{both} ends. Split, say at $2$:
        $\int_{1}^{2}+\int_{2}^{\infty}$. Partial fractions:
        $\dfrac{1}{x(x-1)}=\dfrac{1}{x-1}-\dfrac{1}{x}$, antiderivative
        $\ln\left|\dfrac{x-1}{x}\right|$. Near $x=1^{+}$ this
        $\to-\infty$, so the lower piece \textbf{diverges} and the
        whole integral diverges.
\end{solutions}

\section{Additional Problems}

Decide whether each integral converges; if it does, evaluate it.

\begin{problems}
  \item $\displaystyle\int_{1}^{\infty} \frac{1}{x^{3/2}}\,\mathrm{d}x$
  \item $\displaystyle\int_{1}^{\infty} \frac{1}{x^{4}}\,\mathrm{d}x$
  \item $\displaystyle\int_{2}^{\infty} \mathrm{e}^{-3x}\,\mathrm{d}x$
  \item $\displaystyle\int_{0}^{\infty} \mathrm{e}^{-x/2}\,\mathrm{d}x$
  \item $\displaystyle\int_{0}^{\infty} \frac{1}{(1+x)^{3}}\,\mathrm{d}x$
  \item $\displaystyle\int_{0}^{1} \frac{1}{x^{2/3}}\,\mathrm{d}x$
  \item $\displaystyle\int_{0}^{1} \frac{\ln x}{x}\,\mathrm{d}x$
  \item $\displaystyle\int_{0}^{\infty} t^{3}\,\mathrm{e}^{-t}\,\mathrm{d}t$
  \item $\displaystyle\int_{0}^{\infty} t\,\mathrm{e}^{-2t}\,\mathrm{d}t$
  \item $\displaystyle\int_{-\infty}^{0} x\,\mathrm{e}^{x}\,\mathrm{d}x$
  \item $\displaystyle\int_{\mathrm{e}}^{\infty} \frac{1}{x(\ln x)^{3}}\,\mathrm{d}x$
  \item $\displaystyle\int_{0}^{4} \frac{1}{\sqrt{x}}\,\mathrm{d}x$
  \item $\displaystyle\int_{0}^{9} \frac{1}{\sqrt{9-x}}\,\mathrm{d}x$
  \item $\displaystyle\int_{0}^{1} \frac{\ln x}{\sqrt{x}}\,\mathrm{d}x$
  \item $\displaystyle\int_{0}^{\infty} x^{2}\,\mathrm{e}^{-2x}\,\mathrm{d}x$
  \item $\displaystyle\int_{1}^{\infty} \frac{(\ln x)^{2}}{x^{2}}\,\mathrm{d}x$
  \item $\displaystyle\int_{2}^{\infty} \frac{1}{(x-1)^{2}}\,\mathrm{d}x$
  \item Present value of a perpetuity paying \$2{,}000 per year
        continuously at $r=0.08$.
  \item $\displaystyle\int_{3}^{\infty} \frac{1}{x^{2}-4}\,\mathrm{d}x$
  \item $\displaystyle\int_{0}^{1} \frac{1}{(1-x)^{1/3}}\,\mathrm{d}x$
\end{problems}

\section{Additional Solutions}

\begin{solutions}
  \item $p=3/2>1$, so converges to $\dfrac{1}{p-1}=2$. Check:
        $\displaystyle\bigl[-2x^{-1/2}\bigr]_{1}^{b}=-2/\sqrt{b}+2\to 2$.

  \item $\displaystyle\Bigl[-\frac{1}{3x^{3}}\Bigr]_{1}^{b}
        =-\frac{1}{3b^{3}}+\frac{1}{3}\to\frac{1}{3}$. Converges to
        $1/3$.

  \item $\displaystyle\Bigl[-\frac{1}{3}\mathrm{e}^{-3x}\Bigr]_{2}^{b}
        =-\frac{1}{3}\mathrm{e}^{-3b}+\frac{1}{3}\mathrm{e}^{-6}
        \to\frac{1}{3}\mathrm{e}^{-6}$. Converges to $\mathrm{e}^{-6}/3$.

  \item $\displaystyle\bigl[-2\mathrm{e}^{-x/2}\bigr]_{0}^{b}
        =-2\mathrm{e}^{-b/2}+2\to 2$. Converges to $2$. (Matches $1/r$
        with $r=1/2$.)

  \item Antiderivative $-\dfrac{1}{2(1+x)^{2}}$.
        $\displaystyle\Bigl[-\frac{1}{2(1+x)^{2}}\Bigr]_{0}^{b}
        \to 0+\frac{1}{2}=\frac{1}{2}$. Converges to $1/2$.

  \item $p=2/3<1$, so converges.
        $\displaystyle\int_{t}^{1} x^{-2/3}\,\mathrm{d}x
        =\bigl[3x^{1/3}\bigr]_{t}^{1}=3-3t^{1/3}\to 3$. Converges to
        $3$.

  \item Antiderivative $\dfrac{(\ln x)^{2}}{2}$.
        $\displaystyle\Bigl[\frac{(\ln x)^{2}}{2}\Bigr]_{t}^{1}
        =0-\frac{(\ln t)^{2}}{2}\to-\infty$ as $t\to 0^{+}$. Diverges.

  \item By the gamma identity,
        $\int_{0}^{\infty} t^{n}\mathrm{e}^{-t}\,\mathrm{d}t=n!$. Here
        $n=3$, so the value is $6$.

  \item By the formula
        $\int_{0}^{\infty} t\mathrm{e}^{-rt}\,\mathrm{d}t=1/r^{2}$ with
        $r=2$: $1/4$.

  \item Antiderivative $(x-1)\mathrm{e}^{x}$.
        $\bigl[(x-1)\mathrm{e}^{x}\bigr]_{a}^{0}=-1-(a-1)\mathrm{e}^{a}$.
        As $a\to-\infty$, $(a-1)\mathrm{e}^{a}\to 0$, so the limit is
        $-1$. Converges to $-1$.

  \item $u=\ln x$, $\mathrm{d}u=\mathrm{d}x/x$. Limits $u:1\to\infty$.
        \[
          \int_{1}^{\infty} \frac{1}{u^{3}}\,\mathrm{d}u
            =\Bigl[-\frac{1}{2u^{2}}\Bigr]_{1}^{\infty}=\frac{1}{2} .
        \]
        Converges to $1/2$.

  \item $p=1/2<1$. $\displaystyle\bigl[2\sqrt{x}\bigr]_{t}^{4}
        =4-2\sqrt{t}\to 4$. Converges to $4$.

  \item $u=9-x$, $\mathrm{d}u=-\mathrm{d}x$. Antiderivative of
        $-(9-x)^{-1/2}\cdot(-1)=(9-x)^{-1/2}$ is $-2\sqrt{9-x}$.
        \[
          \bigl[-2\sqrt{9-x}\bigr]_{0}^{t}=-2\sqrt{9-t}+6\to 6 .
        \]
        Converges to $6$ as $t\to 9^{-}$.

  \item By parts: $u=\ln x$, $v'=x^{-1/2}$, $v=2\sqrt{x}$, $u'=1/x$.
        \[
          \int \frac{\ln x}{\sqrt{x}}\,\mathrm{d}x
            =2\sqrt{x}\ln x-\int \frac{2}{\sqrt{x}}\,\mathrm{d}x
            =2\sqrt{x}\ln x-4\sqrt{x}+C .
        \]
        As $t\to 0^{+}$: $\sqrt{t}\ln t\to 0$, $\sqrt{t}\to 0$. At
        $x=1$: $0-4=-4$. So the integral converges to $-4$.

  \item By the formula
        $\int_{0}^{\infty} t^{n}\mathrm{e}^{-rt}\,\mathrm{d}t=n!/r^{n+1}$
        with $n=2$, $r=2$: $2/8=1/4$.

  \item By parts twice. From Section 3 problem 15,
        $\int \frac{\ln x}{x^{2}}\,\mathrm{d}x=-\frac{\ln x+1}{x}$.
        Apply parts to the original: $u=(\ln x)^{2}$, $v'=x^{-2}$,
        $v=-x^{-1}$, $u'=2\ln x/x$.
        \[
          \int \frac{(\ln x)^{2}}{x^{2}}\,\mathrm{d}x
            =-\frac{(\ln x)^{2}}{x}+2\int \frac{\ln x}{x^{2}}\,\mathrm{d}x
            =-\frac{(\ln x)^{2}+2\ln x+2}{x}+C .
        \]
        As $b\to\infty$, $(\ln b)^{k}/b\to 0$ for any $k$. So the value
        is $0-(-2)=2$.

  \item Antiderivative $-\dfrac{1}{x-1}$.
        \[
          \Bigl[-\frac{1}{x-1}\Bigr]_{2}^{b}=-\frac{1}{b-1}+1\to 1 .
        \]
        Converges to $1$.

  \item Perpetuity formula: $PV=D/r=2000/0.08=25{,}000$.

  \item Partial fractions:
        $\dfrac{1}{x^{2}-4}
        =\dfrac{1}{4}\left(\dfrac{1}{x-2}-\dfrac{1}{x+2}\right)$, with
        antiderivative $\dfrac{1}{4}\ln\left|\dfrac{x-2}{x+2}\right|$.
        \[
          \Bigl[\tfrac{1}{4}\ln\tfrac{x-2}{x+2}\Bigr]_{3}^{b}
            =\tfrac{1}{4}\ln\tfrac{b-2}{b+2}-\tfrac{1}{4}\ln\tfrac{1}{5} .
        \]
        As $b\to\infty$, $(b-2)/(b+2)\to 1$, so $\ln\to 0$. Value:
        $-\tfrac{1}{4}\ln\tfrac{1}{5}=\dfrac{\ln 5}{4}$.

  \item $u=1-x$, $\mathrm{d}u=-\mathrm{d}x$. Antiderivative
        $-\tfrac{3}{2}(1-x)^{2/3}$.
        \[
          \Bigl[-\tfrac{3}{2}(1-x)^{2/3}\Bigr]_{0}^{t}
            =-\tfrac{3}{2}(1-t)^{2/3}+\tfrac{3}{2}\to\tfrac{3}{2} .
        \]
        Converges to $3/2$.
\end{solutions}
